\documentclass[11pt]{article}

\usepackage[letterpaper,margin=1in]{geometry}
\usepackage{amsmath}
\usepackage{amssymb}
\usepackage{graphicx}

\graphicspath{{figs/}}

\setlength{\parindent}{1.5em}
\setlength{\parskip}{2pt}

\begin{document}

% =========================================================================
%  Section 2.2.2 -- rewritten as feasibility (preliminary results).
%  Style mirrors 2.2.1 and 2.2.3: narrative prose, embedded numbers,
%  short figure captions, no equations, no sub-headers, no (i)(ii)(iii)
%  bulleted demonstration claims.  Placeholder citations use the
%  [cite: TBD:xxx] convention to be resolved in the bibliography pass.
% =========================================================================

\noindent\textbf{2.2.2 Preliminary Process-Level Digital Twin (RIT)}\par

To complement WPI's experimental demonstration, the RIT team built a
preliminary process-level digital twin of the alkaline-leaching and
HCl-precipitation Si recycling sequence and exercised it against the WPI
endpoint reported in ref.~5. The twin couples a surface-reaction
shrinking-core kinetics model for Si dissolution in NaOH with an
aqueous-equilibrium speciation model for pH-controlled selective
precipitation, and runs in under one second on a workstation, fast
enough to serve as the inner loop of the AI-guided optimization
pipeline proposed in Subtasks~1.3, 2.4, 3.3, and 4.3.

The Si dissolution kinetics model uses a literature activation energy of
60~kJ/mol and a first-order NaOH dependence
[cite: TBD:Si-alkaline-kinetics]. The pre-exponential factor was fixed
by a single calibration to WPI's measured endpoint of 98\,\% Si
conversion at 5~M NaOH, $80\,^\circ$C, 60~min reported in ref.~5. With
no further tuning, the model then predicts conversion trajectories at
other operating points, as shown in Figure~\ref{fig:leach}: 98\,\%
conversion is reached in 99~min at 3~M NaOH and in 42~min at 7~M NaOH,
providing a direct estimate of the trade-off between leach time and
reagent consumption for WPI's upcoming scale-up trials.

\begin{figure}[h]
  \centering
  \includegraphics[width=0.58\textwidth]{leaching_kinetics.pdf}
  \caption{Predicted Si conversion vs leaching time at $80\,^\circ$C for
    three NaOH concentrations. Open marker: WPI calibration point (5~M,
    60~min, ref.~5).\label{fig:leach}}
\end{figure}

The selective-precipitation model computes equilibrium solubilities of
Si(IV) and Al(III) above amorphous SiO$_{2}$ and Al(OH)$_{3}$ using
textbook $\log K$ values from the MINTEQ database
[cite: TBD:MINTEQ]. As shown in Figure~\ref{fig:select}, the model
identifies a sharp selectivity window at $\mathrm{pH} < 2$ in which Si
is forced to the SiO$_{2}$ saturation limit of $\sim$2~mM
(near-quantitative precipitation) while Al remains many orders of
magnitude undersaturated. Operating below $\mathrm{pH}=1$ widens the
Si/Al saturation contrast by another five orders of magnitude,
suppressing Al co-precipitation. This is consistent with WPI's
empirically chosen acidification target of $\mathrm{pH} < 1$ and
explains the observed 6N SiO$_{2}$ purity with Al, Pb, and Sn retained
in the leachate (\S2.2.1).

\begin{figure}[h]
  \centering
  \includegraphics[width=0.58\textwidth]{selective_precipitation.pdf}
  \caption{Predicted equilibrium solubilities of Si(IV) and Al(III) vs
    pH. Shaded band: selective SiO$_2$ recovery window
    ($\mathrm{pH} < 2$).\label{fig:select}}
\end{figure}

These preliminary results show that the WPI process chemistry is
captured by tractable continuum-thermodynamic modules using literature
parameters and a single calibration point, and that the same modular
structure extends directly to the thin-film PV chemistries proposed in
Subtask~2.2 (BU) by swapping the species library and $\log K$ values.
The twin is implemented on top of the open-source MOOSE-based
\textsc{FEECM} framework [cite: TBD:FEECM-repo] developed by Co-PI Tu's
group for coupled electrochemical transport and reaction-kinetics
modeling, and will serve as the foundation for the reactor-scale,
AI-augmented digital twin proposed in BP2 and BP3, running on CewLab GPU
compute, RIT's SPORC HPC platform, the New York State Empire AI
allocation, and ACCESS/NSF cyberinfrastructure.

% =========================================================================
% Appended drafts for the remaining highlighted (XXX) parts.
% Each block below carries a "[Insert at ...]" header so the user can drop
% it into the matching location in the Word source.
% =========================================================================

\vspace{2em}\noindent\rule{\textwidth}{1pt}

\begin{center}\textsc{Drafts for the remaining highlighted (XXX) sections}\end{center}

\noindent The blocks below are drafts for each XXX placeholder
identified in the source draft. Each carries a header indicating where
it belongs in the Word source.

\vspace{1em}\noindent\rule{0.4\textwidth}{0.4pt}\vspace{0.6em}

% -------------------------------------------------------------------------
% Subtask 1.3  (under Task 1.0, BP1)
% -------------------------------------------------------------------------

\noindent\textit{[Insert at \S3.5 Task 1.0 (BP1), Subtask 1.3]}\par

\noindent\textbf{Subtask 1.3:} Development of a digital twin framework
for Si-based photovoltaic recycling, including reaction kinetics, mass
transfer, and process optimization under continuous processing
conditions. (RIT)\par

\noindent\textbf{Subtask Summary:} The RIT team will develop a
calibrated reactor-scale digital twin of the Si-based PV recycling
sequence (NaOH leaching $\rightarrow$ HCl precipitation), extending the
preliminary process-level twin demonstrated in \S2.2.2 to include solids
loading, mixing, residence-time distribution, temperature uniformity,
and pH trajectory at scaled processing conditions. The twin couples a
multi-species mass-balance and pH-speciation core with a reactor-scale
transport model and a physics-constrained AI surrogate trained on
simulation outputs to recommend operating windows that maximize Si
recovery and 6N purity while minimizing NaOH and HCl consumption. The
model will be calibrated against WPI's ICP, pH, conductivity, and
yield/purity data collected in Subtasks~1.1 and~1.2 and is implemented
on the open-source MOOSE-based \textsc{FEECM} framework developed by
Co-PI Tu's group.\par

\noindent\textbf{Milestone 1.3.1:} Deliver a calibrated reactor-scale
digital twin for Si-based PV recycling that predicts Si conversion, 6N
SiO$_2$ yield, and reagent consumption within $\pm$5\,\% of WPI's
scaled-batch measurements across the operating range
$[\mathrm{NaOH}] = 3$--$7$\,M, $T = 60$--$90\,^\circ$C, precipitation
$\mathrm{pH} < 2$.\par

\noindent\textbf{Milestone 1.3.2:} Deliver an AI-surrogate-driven
operating-window recommendation that achieves the targeted $>$95\,\% Si
recovery and 6N purity at WPI's scaled processing conditions in a
single recommended setting, with the recommendation accompanied by
quantified uncertainty bounds.

\vspace{1em}\noindent\rule{0.4\textwidth}{0.4pt}\vspace{0.6em}

% -------------------------------------------------------------------------
% Subtask 2.4  (under Task 2.0, BP1)
% -------------------------------------------------------------------------

\noindent\textit{[Insert at \S3.5 Task 2.0 (BP1), Subtask 2.4]}\par

\noindent\textbf{Subtask 2.4:} Development of a digital twin framework
for thin-film photovoltaic recycling to model selective recovery
behavior, impurity interactions, and process stability for critical
material recovery. (RIT)\par

\noindent\textbf{Subtask Summary:} The RIT team will extend the
Si-based digital twin developed in Subtask~1.3 to the thin-film PV
chemistries investigated by the BU team in Subtask~2.2, with the goal
of predicting selective recovery of critical materials (Ga, In, Ge,
Cu) from CIGS and related thin-film waste streams. The extension
augments the species library with the new soluble and solid phases,
incorporates the multi-ion interactions and competing co-precipitation
pathways that arise in mixed thin-film leachates, and adds an off-gas
equilibrium submodel that predicts H$_2$Se generation under acidic
conditions (linked to Subtask~2.3). The model will be calibrated
against BU's ICP, XRD, SEM-EDS, and XPS feedstock and product
datasets collected in Subtask~2.1.\par

\noindent\textbf{Milestone 2.4.1:} Deliver a calibrated digital twin
for thin-film PV recycling that predicts $>$98\,\% recovery efficiency
and $>$4N purity for Ga, In, Ge, and Cu across representative CIGS
feedstock compositions, validated against BU scaled-batch
measurements within $\pm$5\,\%.\par

\noindent\textbf{Milestone 2.4.2:} Deliver a model-based safe-operating
window that predicts H$_2$Se generation rates as a function of pH, acid
addition rate, temperature, and feedstock Se content, with confirmed
agreement against BU's off-gas measurements (Subtask~2.3).

\vspace{1em}\noindent\rule{0.4\textwidth}{0.4pt}\vspace{0.6em}

% -------------------------------------------------------------------------
% Subtask 3.3  (under Task 3.0, BP2)
% -------------------------------------------------------------------------

\noindent\textit{[Insert at \S3.5 Task 3.0 (BP2), Subtask 3.3]}\par

\noindent\textbf{Subtask 3.3:} Development of a digital twin platform
incorporating continuum modeling, reaction kinetics, and AI-assisted
optimization. (RIT)\par

\noindent\textbf{Subtask Summary:} The RIT team will integrate the
Si-based (Subtask~1.3) and thin-film (Subtask~2.4) digital twins into a
unified, chemistry-adaptable PV recycling platform capable of handling
mixed Si-based and thin-film feedstocks. The platform exposes the
elemental mass-balance, reaction-kinetics, reactor-transport, and
AI-surrogate optimization layers as composable modules, so that a
single user input--feedstock composition, operating conditions, target
purity--drives the full simulation-to-recommendation workflow. The
physics-constrained AI surrogate will recommend mixed-feedstock
operating conditions that maximize recovery and purity while minimizing
chemical consumption, embodied energy, and waste generation, and will
be calibrated against the 1\,kg mixed-feedstock validation data from
Subtask~3.4. The platform will be made available as a shared
computational resource accessible to all project teams.\par

\noindent\textbf{Milestone 3.3.1:} Deliver a validated digital twin
platform integrating the Si-based and thin-film recycling models that
predicts mixed-feedstock recovery efficiency and product purity within
$\pm$5\,\% of the 1\,kg measurement endpoints from Subtask~3.4.\par

\noindent\textbf{Milestone 3.3.2:} Deliver an AI-assisted optimization
interface that reduces the number of experimental iterations required
to identify scalable operating conditions for a new feedstock blend by
at least 50\,\% relative to a Latin-hypercube sweep of equivalent
coverage.

\vspace{1em}\noindent\rule{0.4\textwidth}{0.4pt}\vspace{0.6em}

% -------------------------------------------------------------------------
% Subtask 4.3  (under Task 4.0, BP3)
% -------------------------------------------------------------------------

\noindent\textit{[Insert at \S3.5 Task 4.0 (BP3), Subtask 4.3]}\par

\noindent\textbf{Subtask 4.3:} Computational modeling of the pilot-scale
PV recycling. (RIT)\par

\noindent\textbf{Subtask Summary:} The RIT team will apply the
validated digital twin platform from Subtask~3.3 to model the 5\,kg and
pilot-scale PV recycling reactors developed in Subtasks~4.1, 4.2,
and~4.4. Modeling at this scale captures the scale-up--sensitive
phenomena that are absent at beaker scale: reactor geometry effects on
mixing, heat-transfer--limited temperature uniformity, residence-time
distributions in continuous operation, and impurity carryover under
increased solids loading. The model will be used to guide reactor
design and operating-parameter selection prior to the BU/WPI/industry
5\,kg validation runs, to predict throughput, energy demand, and reagent
consumption as inputs to the TEA/LCA (Task~5.0), and to identify safe
operating envelopes for H$_2$Se emission control under scaled
conditions.\par

\noindent\textbf{Milestone 4.3.1:} Deliver scale-up predictions for the
5\,kg and pilot-scale PV recycling processes--recovery efficiency,
product purity, throughput, energy demand, reagent consumption, and
H$_2$Se emission risk--validated against industry-partner pilot-scale
measurements within $\pm$10\,\%.\par

\noindent\textbf{Milestone 4.3.2:} Deliver model-based operating
recommendations and uncertainty bounds for safe and economically viable
continuous operation at pilot scale, used as direct input to the TEA/LCA
in Task~5.0.

\vspace{1em}\noindent\rule{0.4\textwidth}{0.4pt}\vspace{0.6em}

% -------------------------------------------------------------------------
% Section 4.1 -- Dr. Tu bio
% -------------------------------------------------------------------------

\noindent\textit{[Insert at \S4.1 Team's Unique Qualifications and Expertise, after Prof.\ Yan Wang's bio]}\par

\noindent Dr.\ Qingsong (Howard) Tu (Co-PI) is Assistant Professor of
Mechanical Engineering at Rochester Institute of Technology (RIT). His
research focuses on physics-informed digital twins, continuum-scale
transport modeling, and machine-learning--assisted process optimization
for electrochemical and materials-manufacturing systems, with
applications spanning battery and electrolyte recycling, lithium-ion
and solid-state batteries, electrochemical ion insertion, and
sustainable manufacturing of advanced materials. He has published more
than 40 peer-reviewed papers and applied for 5 patents, with over
3,500 citations. Dr.\ Tu directs the Clean Energy and Water Lab
(CewLab) at RIT, which combines experimental and computational
facilities for the coupled study of electro-chemo-mechanics,
electrochemical ion transport, and process scale-up. In the proposed
project, Dr.\ Tu leads the digital twin development, AI-assisted
process optimization, and computational scale-up tasks
(Subtasks~1.3, 2.4, 3.3, and~4.3).

\vspace{1em}\noindent\rule{0.4\textwidth}{0.4pt}\vspace{0.6em}

% -------------------------------------------------------------------------
% Section 4.2 -- CewLab compute completion
% -------------------------------------------------------------------------

\noindent\textit{[Insert at \S4.2 Team's existing Equipment and Facilities, completing the trailing sentence ``The following are computational resources\ldots'']}\par

\noindent Local GPU workstations and a small CPU cluster in CewLab
support model development, parameter sweeps, and small-batch surrogate
training. RIT's SPORC (Sustainable Power Operations and Research
Center) high-performance computing platform provides multi-node CPU
and GPU resources for production-scale multiphysics simulations and
reactor-scale digital twin runs. Co-PI Tu's group holds an allocation
on the New York State Empire AI consortium, which provides access to
large-scale GPU compute for training the physics-constrained AI
surrogates that drive the operating-window recommendations in this
project. National cyberinfrastructure allocations through the NSF
ACCESS program provide additional capacity for parallel runs that
exceed local capacity. Together these resources form a tiered compute
stack--development on CewLab GPUs, production on SPORC, and large-scale
AI training on Empire AI/ACCESS--that scales with the simulation and
surrogate-training workload of the proposed project.

\vspace{1em}\noindent\rule{0.4\textwidth}{0.4pt}\vspace{0.6em}

% -------------------------------------------------------------------------
% Section 4.4 -- Tu time commitment
% -------------------------------------------------------------------------

\noindent\textit{[Insert at \S4.4 Time Commitment of the Key Team Members, completing the trailing ``Dr.\ Tu\ldots'']}\par

\noindent Dr.\ Tu is the Co-PI for the project and will spend
\textbf{1~month} per year on the project.\footnote{Placeholder value
matching the typical Co-PI commitment on DOE awards; please confirm
against the negotiated effort.} In addition, Dr.\ Tu will advise
\textbf{1~Ph.D.\ student} working on this project, who will work 50\,\%
time during semesters and 100\,\% effort during summer.

\end{document}
